% Part: computability
% Chapter: recursive-functions
% Section: notations-pr-functions

\documentclass[../../../include/open-logic-section]{subfiles}

\begin{document}

\olfileid[pt]{cmp}{rec}{not}
\olsection{Notações para funções recursivas primitivas}

Uma vantagem de dispor da descrição indutiva precisa das funções recursivas
primitivas é podermos descrevê-las sistematicamente. Por exemplo, podemos
atribuir uma ``notação'' a cada uma dessas funções do modo seguinte. Usamos os
símbolos $\Zero$, $\Succ$ e $\Proj{n}{i}$ para a função zero, a função sucessor
e as projeções. Suponha agora que $h$ seja definida por composição a partir de
uma função de $k$ argumentos~$f$ e de funções de $n$ argumentos $g_0$,
\dots,~$g_{k-1}$, e que tenhamos atribuído às últimas funções as notações $F$, $G_0$,
\dots,~$G_{k-1}$. Usando então um novo símbolo $\fn{Comp}_{k,n}$, podemos
denotar a função $h$ por $\fn{Comp}_{k,n}[F,G_0,\dots,G_{k-1}]$.

Para funções definidas por recursão primitiva, podemos usar notações análogas.
Suponha que a função $(k+1)$-ária~$h$ seja definida por recursão primitiva a
partir da função $k$-ária~$f$ e da função $(k+2)$-ária~$g$, e que as notações
atribuídas a $f$ e~$g$ sejam, respectivamente, $F$ e~$G$. A notação atribuída
a~$h$ é então $\fn{Rec}_k[F,G]$.

Lembre-se de que a função adição é definida por recursão primitiva como
\begin{align*}
  \Add(x_0, 0) & = \Proj{1}{0}(x_0) = x_0\\
  \Add(x_0, y+1) & = \Succ(\Proj{3}{2}(x_0, y, \Add(x_0, y))) = \Add(x_0, y) +1
\end{align*}
Aqui o papel de~$f$ é desempenhado por $\Proj{1}{0}$, e o papel de~$g$, por
$\Succ(\Proj{3}{2}(x_0, y, z))$. A esta última função atribui-se a notação
$\fn{Comp}_{1,3}[\Succ,\Proj{3}{2}]$, pois ela resulta da definição de uma
função por composição a partir da função $1$-ária~$\Succ$ e da função
$3$-ária~$\Proj{3}{2}$. Com essa convenção, podemos denotar a função adição por
\[
\fn{Rec}_1[\Proj{1}{0},\fn{Comp}_{1,3}[\Succ,\Proj{3}{2}]].
\]
Dispor dessas notações às vezes é útil, p.ex., ao enumerar as funções
recursivas primitivas.

\begin{prob}
  Dê a notação recursiva primitiva completa para~$\Mult$.
\end{prob}

\end{document}
